\documentclass[12pt,letterpaper]{article}

% -------------------------------------------------
% Formato general tipo APA 7.ª edición
% -------------------------------------------------
\usepackage[utf8]{inputenc}
\usepackage[T1]{fontenc}
\usepackage[spanish,es-tabla]{babel}
\usepackage{geometry}
\usepackage{setspace}
\usepackage{indentfirst}
\usepackage{mathptmx}
\usepackage{xcolor}
\usepackage{hyperref}
\usepackage{longtable}
\usepackage{array}
\usepackage{booktabs}
\usepackage{enumitem}
\usepackage{fancyhdr}
\usepackage{titlesec}
\usepackage{listings}
\usepackage[most]{tcolorbox}
\usepackage{graphicx}
\usepackage{ragged2e}
\usepackage{seqsplit}
\usepackage{float}
\usepackage{etoolbox}

% Márgenes APA: 1 pulgada
\geometry{margin=1in}
\setlength{\headheight}{15pt}

% Interlineado APA: doble espacio
\doublespacing

% Sangría APA: primera línea de cada párrafo 0.5 pulgadas
\setlength{\parindent}{0.5in}
\setlength{\parskip}{0pt}

% Alineación recomendada por APA: margen izquierdo alineado y derecho irregular
\raggedright

% Encabezado y numeración APA
\pagestyle{fancy}
\fancyhf{}
\rhead{\thepage}
\renewcommand{\headrulewidth}{0pt}

% Colores y enlaces
\definecolor{apaBlue}{HTML}{1F4E79}
\definecolor{codeBg}{HTML}{F6F8FA}
\definecolor{codeFrame}{HTML}{D0D7DE}
\definecolor{quoteBg}{HTML}{F7F7F7}
\definecolor{quoteFrame}{HTML}{BDBDBD}

\hypersetup{
    colorlinks=true,
    linkcolor=black,
    urlcolor=apaBlue,
    citecolor=black,
    pdftitle={Diccionario de Datos: Jeiggar Vacation},
    pdfauthor={NICK ORTEGA}
}

% Jerarquía de títulos compatible con APA
\titleformat{\section}
  {\normalfont\centering\bfseries\Large}
  {\thesection.}{0.5em}{}

\titleformat{\subsection}
  {\normalfont\bfseries\large}
  {\thesubsection}{0.5em}{}

\titleformat{\subsubsection}
  {\normalfont\bfseries\itshape\normalsize}
  {\thesubsubsection}{0.5em}{}

\titlespacing*{\section}{0pt}{1.5em}{1em}
\titlespacing*{\subsection}{0pt}{1.2em}{0.8em}
\titlespacing*{\subsubsection}{0pt}{1em}{0.6em}

% Cada capítulo/sección principal inicia en una nueva hoja
\pretocmd{\section}{\clearpage}{}{}

% Código fuente
\lstdefinelanguage{sql}{
    keywords={SELECT,FROM,WHERE,AND,OR,TRUE,FALSE,NULL,CREATE,TABLE,FUNCTION,RETURNS,TRIGGER,LANGUAGE,STABLE,INSERT,UPDATE,DELETE,PRIMARY,KEY,FOREIGN,REFERENCES,CHECK,UNIQUE,DEFAULT},
    sensitive=false,
    comment=[l]{--},
    morestring=[b]'
}

\lstdefinestyle{codeblock}{
    backgroundcolor=\color{codeBg},
    frame=single,
    rulecolor=\color{codeFrame},
    basicstyle=\ttfamily\small,
    breaklines=true,
    breakatwhitespace=false,
    columns=fullflexible,
    keepspaces=true,
    showstringspaces=false,
    tabsize=2
}
\lstset{style=codeblock}

\newtcolorbox{infobox}{
    colback=quoteBg,
    colframe=quoteFrame,
    boxrule=0.8pt,
    arc=1mm,
    left=2mm,
    right=2mm,
    top=1mm,
    bottom=1mm
}

% Texto técnico con posibilidad de corte de línea dentro de tablas
\DeclareRobustCommand{\code}[1]{\texorpdfstring{\begingroup\ttfamily\footnotesize\def\_{\textunderscore\allowbreak}#1\endgroup}{#1}}
\newcolumntype{L}[1]{>{\RaggedRight\arraybackslash}p{#1}}
\newcommand{\dash}{---}
\newcommand{\yes}{Sí}
\newcommand{\no}{No}

% Ajuste para tablas largas
\setlength{\LTpre}{0.8em}
\setlength{\LTpost}{0.8em}
\renewcommand{\arraystretch}{1.25}
\BeforeBeginEnvironment{longtable}{\begingroup\singlespacing\footnotesize\setlength{\tabcolsep}{3pt}}
\AfterEndEnvironment{longtable}{\endgroup}

\begin{document}

% -------------------------------------------------
% Portada tipo APA 7 estudiante
% -------------------------------------------------
\begin{titlepage}
    \centering
    \vspace*{1.5in}

    {\bfseries\Large Manual de Código: Plataforma Web de Gestión de Destinos y Administración de Servicios Turísticos\par}
    \vspace{0.5in}

    {Estudiantes de IV Semestre\par}
    \vspace{0.2in}

    {Fundación de Estudios Superiores Comfanorte (FESC)\par}
    {Ingeniería de Software\par}
    \vspace{0.2in}

    {Diccionario de Datos\par}
    \vspace{0.2in}

    {Docente: Norly Alexandra Aguilar Quintero \par}
    \vspace{0.2in}

    {Mayo de 2026\par}
\end{titlepage}

\newpage
\tableofcontents
\newpage

\section{Introducción}

La base de datos de \textbf{Jeiggar Vacation} está alojada en Supabase, sobre PostgreSQL 15 o superior, y soporta una plataforma de turismo que permite explorar destinos nacionales e internacionales a través de un mapa interactivo. El sistema gestiona una jerarquía geográfica de tres niveles: países, ciudades y destinos.

\begin{longtable}{L{0.45\textwidth} L{0.35\textwidth}}
\toprule
\textbf{Elemento} & \textbf{Cantidad} \\
\midrule
Tablas & 4 \\
Funciones SQL & 2 \\
Triggers & 3 \\
Índices & 9 \\
Políticas RLS & 14 \\
Extensiones & 4: \code{pg\_stat\_statements}, \code{pgcrypto}, \code{supabase\_vault}, \code{uuid-ossp} \\
\bottomrule
\end{longtable}

\section{Modelo Relacional}

El esquema sigue una jerarquía geográfica donde cada país contiene ciudades y cada ciudad contiene destinos. Los destinos pueden también relacionarse directamente con un país para representar destinos internacionales sin ciudad específica. La tabla \code{profiles} extiende a los usuarios de Supabase Auth con el rol de administrador.

\begin{longtable}{L{0.22\textwidth} L{0.22\textwidth} L{0.18\textwidth} L{0.28\textwidth}}
\toprule
\textbf{Tabla origen} & \textbf{Tabla destino} & \textbf{Cardinalidad} & \textbf{FK / Detalle} \\
\midrule
\code{countries} & \code{cities} & 1:N & \code{countries.id} $\rightarrow$ \code{cities.country\_id}; eliminación en cascada. \\
\code{countries} & \code{destinations} & 1:N & \code{countries.id} $\rightarrow$ \code{destinations.country\_id}. \\
\code{cities} & \code{destinations} & 1:N & \code{cities.id} $\rightarrow$ \code{destinations.city\_id}; eliminación en cascada. \\
\code{auth.users} & \code{profiles} & 1:1 & \code{auth.users.id} $\rightarrow$ \code{profiles.id}; eliminación en cascada. \\
\bottomrule
\end{longtable}

\section{Descripción de Tablas}

\subsection{\code{countries}: Países}

Almacena los países disponibles en la plataforma. Es el nivel raíz de la jerarquía geográfica.

\begin{longtable}{L{0.19\textwidth} L{0.20\textwidth} L{0.09\textwidth} L{0.17\textwidth} L{0.25\textwidth}}
\toprule
\textbf{Campo} & \textbf{Tipo} & \textbf{Nulo} & \textbf{Default} & \textbf{Descripción} \\
\midrule
\code{id} & \code{uuid} & NO & \code{gen\_random\_uuid()} & Identificador único del país; llave primaria. \\
\code{name} & \code{text} & NO & \dash & Nombre oficial del país. Ejemplo: Colombia. \\
\code{slug} & \code{text} & NO & \dash & Identificador URL amigable único. Ejemplo: \code{colombia}. \\
\code{lat} & \code{double precision} & NO & \dash & Latitud del centro geográfico del país. \\
\code{lng} & \code{double precision} & NO & \dash & Longitud del centro geográfico del país. \\
\code{zoom} & \code{numeric(4,2)} & NO & \dash & Nivel de zoom inicial en MapLibre/Mapbox al centrar la vista en el país. Rango típico: 4--7. \\
\code{image\_url} & \code{text} & SÍ & \code{NULL} & URL de imagen representativa del país. \\
\code{description} & \code{text} & SÍ & \code{NULL} & Descripción general del país para mostrar en la interfaz. \\
\code{is\_active} & \code{boolean} & NO & \code{true} & Controla la visibilidad pública. Solo los registros activos son accesibles sin autenticación admin. \\
\code{sort\_order} & \code{integer} & NO & \code{0} & Orden de presentación en listados. Menor número equivale a mayor prioridad. \\
\code{created\_at} & \code{timestamptz} & NO & \code{now()} & Fecha y hora de creación del registro. \\
\code{updated\_at} & \code{timestamptz} & NO & \code{now()} & Fecha y hora de la última actualización, gestionada automáticamente por \code{trg\_countries\_updated\_at}. \\
\bottomrule
\end{longtable}

\subsubsection*{Restricciones}
\begin{longtable}{L{0.34\textwidth} L{0.22\textwidth} L{0.34\textwidth}}
\toprule
\textbf{Nombre} & \textbf{Tipo} & \textbf{Detalle} \\
\midrule
\code{countries\_pkey} & PRIMARY KEY & \code{id} \\
\code{countries\_slug\_key} & UNIQUE & \code{slug} \\
\bottomrule
\end{longtable}

\subsection{\code{cities}: Ciudades}

Almacena las ciudades o regiones turísticas dentro de cada país.

\begin{longtable}{L{0.19\textwidth} L{0.20\textwidth} L{0.09\textwidth} L{0.17\textwidth} L{0.25\textwidth}}
\toprule
\textbf{Campo} & \textbf{Tipo} & \textbf{Nulo} & \textbf{Default} & \textbf{Descripción} \\
\midrule
\code{id} & \code{uuid} & NO & \code{gen\_random\_uuid()} & Identificador único de la ciudad; llave primaria. \\
\code{country\_id} & \code{uuid} & NO & \dash & Llave foránea hacia \code{countries.id}. Al eliminar el país, las ciudades se eliminan en cascada. \\
\code{name} & \code{text} & NO & \dash & Nombre de la ciudad o región. Ejemplo: Cúcuta. \\
\code{slug} & \code{text} & NO & \dash & Identificador URL amigable. Único por país mediante \code{country\_id + slug}. \\
\code{lat} & \code{double precision} & NO & \dash & Latitud del centro de la ciudad. \\
\code{lng} & \code{double precision} & NO & \dash & Longitud del centro de la ciudad. \\
\code{zoom} & \code{numeric(4,2)} & NO & \dash & Nivel de zoom en el mapa MapLibre al centrar la vista en la ciudad. Rango típico: 10--13. \\
\code{description} & \code{text} & NO & \dash & Descripción de la ciudad para la tarjeta informativa en el mapa. \\
\code{image\_url} & \code{text} & SÍ & \code{NULL} & URL de imagen representativa de la ciudad. \\
\code{is\_active} & \code{boolean} & NO & \code{true} & Visibilidad pública del registro. \\
\code{sort\_order} & \code{integer} & NO & \code{0} & Orden de presentación en listados. \\
\code{created\_at} & \code{timestamptz} & NO & \code{now()} & Fecha de creación. \\
\code{updated\_at} & \code{timestamptz} & NO & \code{now()} & Fecha de última actualización mediante trigger automático. \\
\bottomrule
\end{longtable}

\subsubsection*{Restricciones}
\begin{longtable}{L{0.34\textwidth} L{0.22\textwidth} L{0.34\textwidth}}
\toprule
\textbf{Nombre} & \textbf{Tipo} & \textbf{Detalle} \\
\midrule
\code{cities\_pkey} & PRIMARY KEY & \code{id} \\
\code{cities\_country\_id\_slug\_key} & UNIQUE & \code{country\_id, slug} \\
\code{cities\_country\_id\_fkey} & FOREIGN KEY & \code{country\_id} $\rightarrow$ \code{countries.id} ON DELETE CASCADE \\
\bottomrule
\end{longtable}

\subsection{\code{destinations}: Destinos turísticos}

Es la tabla central de la plataforma. Contiene cada punto de interés turístico con su información geográfica, descriptiva y multimedia. Un destino puede ser nacional, asociado a una ciudad, o internacional, asociado directamente a un país.

\begin{longtable}{L{0.18\textwidth} L{0.20\textwidth} L{0.08\textwidth} L{0.16\textwidth} L{0.28\textwidth}}
\toprule
\textbf{Campo} & \textbf{Tipo} & \textbf{Nulo} & \textbf{Default} & \textbf{Descripción} \\
\midrule
\code{id} & \code{uuid} & NO & \code{gen\_random\_uuid()} & Identificador único del destino; llave primaria. \\
\code{city\_id} & \code{uuid} & SÍ & \code{NULL} & Llave foránea hacia \code{cities.id}. Nulo para destinos internacionales. Cascada al eliminar ciudad. \\
\code{country\_id} & \code{uuid} & SÍ & \code{NULL} & Llave foránea hacia \code{countries.id}. Permite asociar destinos directamente a un país. \\
\code{name} & \code{text} & NO & \dash & Nombre del destino. Ejemplo: Lago de Maracaibo. \\
\code{slug} & \code{text} & NO & \dash & Identificador URL amigable único global. \\
\code{short\_description} & \code{text} & NO & \dash & Descripción corta para tarjetas y previsualizaciones. Máximo recomendado: 160 caracteres. \\
\code{description} & \code{text} & NO & \dash & Descripción completa del destino para la vista de detalle. \\
\code{category} & \code{text} & NO & \dash & Categoría turística. Ejemplos: Naturaleza, Cultura, Aventura, Gastronomía. \\
\code{destination\_type} & \code{text} & NO & \dash & Tipo de destino: \code{national} o \code{international}. Tiene restricción CHECK. \\
\code{lat} & \code{double precision} & NO & \dash & Latitud exacta del destino para el marcador en el mapa. \\
\code{lng} & \code{double precision} & NO & \dash & Longitud exacta del destino para el marcador en el mapa. \\
\code{climate} & \code{text} & SÍ & \code{NULL} & Descripción del clima del destino. Ejemplos: Tropical, frío de montaña. \\
\code{image\_url} & \code{text} & SÍ & \code{NULL} & URL de imagen principal del destino. \\
\code{activities} & \code{text[]} & NO & \code{\{\}} & Arreglo de actividades disponibles. Ejemplos: senderismo, kayak, fotografía. \\
\code{is\_active} & \code{boolean} & NO & \code{true} & Visibilidad pública. \\
\code{is\_featured} & \code{boolean} & NO & \code{false} & Indica si el destino aparece en la sección de destacados en la página de inicio. \\
\code{sort\_order} & \code{integer} & NO & \code{0} & Orden dentro de la ciudad o país. \\
\code{sort\_order\_home} & \code{integer} & NO & \code{0} & Orden en la sección de destacados de la página de inicio. \\
\code{created\_at} & \code{timestamptz} & NO & \code{now()} & Fecha de creación. \\
\code{updated\_at} & \code{timestamptz} & NO & \code{now()} & Fecha de última actualización mediante trigger automático. \\
\bottomrule
\end{longtable}

\subsubsection*{Restricciones}
\begin{longtable}{L{0.34\textwidth} L{0.22\textwidth} L{0.34\textwidth}}
\toprule
\textbf{Nombre} & \textbf{Tipo} & \textbf{Detalle} \\
\midrule
\code{destinations\_pkey} & PRIMARY KEY & \code{id} \\
\code{destinations\_slug\_key} & UNIQUE & \code{slug} \\
\code{destinations\_destination\_type\_check} & CHECK & \code{destination\_type IN ('national', 'international')} \\
\code{destinations\_city\_id\_fkey} & FOREIGN KEY & \code{city\_id} $\rightarrow$ \code{cities.id} ON DELETE CASCADE \\
\code{destinations\_country\_id\_fkey} & FOREIGN KEY & \code{country\_id} $\rightarrow$ \code{countries.id} \\
\bottomrule
\end{longtable}

\subsection{\code{profiles}: Perfiles de usuario}

Extiende la tabla \code{auth.users} de Supabase para asignar roles dentro de la plataforma. Actualmente solo existe el rol \code{admin}, que habilita operaciones de escritura sobre todas las tablas.

\begin{longtable}{L{0.19\textwidth} L{0.20\textwidth} L{0.09\textwidth} L{0.17\textwidth} L{0.25\textwidth}}
\toprule
\textbf{Campo} & \textbf{Tipo} & \textbf{Nulo} & \textbf{Default} & \textbf{Descripción} \\
\midrule
\code{id} & \code{uuid} & NO & \dash & Llave foránea hacia \code{auth.users.id}. También es la llave primaria. Se elimina en cascada al borrar el usuario. \\
\code{role} & \code{text} & NO & \dash & Rol del usuario. Único valor permitido: \code{admin}, mediante restricción CHECK. \\
\code{created\_at} & \code{timestamptz} & SÍ & \code{now()} & Fecha de asignación del perfil. \\
\bottomrule
\end{longtable}

\subsubsection*{Restricciones}
\begin{longtable}{L{0.34\textwidth} L{0.22\textwidth} L{0.34\textwidth}}
\toprule
\textbf{Nombre} & \textbf{Tipo} & \textbf{Detalle} \\
\midrule
\code{profiles\_pkey} & PRIMARY KEY & \code{id} \\
\code{profiles\_role\_check} & CHECK & \code{role = 'admin'} \\
\code{profiles\_id\_fkey} & FOREIGN KEY & \code{id} $\rightarrow$ \code{auth.users.id} ON DELETE CASCADE \\
\bottomrule
\end{longtable}

\section{Índices}

Todos los índices utilizan el método B-Tree y optimizan las consultas más frecuentes de la plataforma.

\begin{longtable}{L{0.30\textwidth} L{0.20\textwidth} L{0.23\textwidth} L{0.17\textwidth}}
\toprule
\textbf{Nombre} & \textbf{Tabla} & \textbf{Columna(s)} & \textbf{Propósito} \\
\midrule
\code{idx\_cities\_country\_id} & \code{cities} & \code{country\_id} & Acelera JOIN y filtros por país. \\
\code{idx\_cities\_slug} & \code{cities} & \code{slug} & Búsqueda rápida por slug de ciudad. \\
\code{idx\_destinations\_city\_id} & \code{destinations} & \code{city\_id} & Acelera JOIN ciudad-destino. \\
\code{idx\_destinations\_category} & \code{destinations} & \code{category} & Filtro por categoría turística. \\
\code{idx\_destinations\_slug} & \code{destinations} & \code{slug} & Búsqueda rápida por slug. \\
\code{idx\_dest\_active} & \code{destinations} & \code{is\_active} & Filtro de registros activos. \\
\code{idx\_dest\_featured} & \code{destinations} & \code{is\_featured} & Filtro de destacados en home. \\
\code{idx\_dest\_type} & \code{destinations} & \code{destination\_type} & Filtro nacional / internacional. \\
\bottomrule
\end{longtable}

\section{Triggers}

\begin{longtable}{L{0.32\textwidth} L{0.18\textwidth} L{0.18\textwidth} L{0.22\textwidth}}
\toprule
\textbf{Nombre} & \textbf{Tabla} & \textbf{Evento} & \textbf{Descripción} \\
\midrule
\code{trg\_countries\_updated\_at} & \code{countries} & BEFORE UPDATE & Llama a \code{set\_updated\_at()} para actualizar automáticamente la columna \code{updated\_at}. \\
\code{trg\_cities\_updated\_at} & \code{cities} & BEFORE UPDATE & Llama a \code{set\_updated\_at()} para actualizar automáticamente la columna \code{updated\_at}. \\
\code{trg\_destinations\_updated\_at} & \code{destinations} & BEFORE UPDATE & Llama a \code{set\_updated\_at()} para actualizar automáticamente la columna \code{updated\_at}. \\
\bottomrule
\end{longtable}

\section{Funciones Almacenadas}

\subsection{\code{set\_updated\_at()}}

\begin{infobox}
\textbf{Tipo:} TRIGGER FUNCTION. \textbf{Lenguaje:} PL/pgSQL. \textbf{Retorna:} \code{trigger}.
\end{infobox}

Asigna el timestamp actual, mediante \code{now()}, al campo \code{updated\_at} del registro modificado. Es invocada por los tres triggers de auditoría.

\subsection{\code{get\_country\_map\_data(p\_country\_slug text)}}

\begin{infobox}
\textbf{Tipo:} SQL FUNCTION. \textbf{Lenguaje:} SQL. \textbf{Volatilidad:} STABLE. \textbf{Retorna:} \code{jsonb}.
\end{infobox}

Es la RPC principal del mapa interactivo. Recibe el slug de un país y devuelve un objeto \code{jsonb} completamente anidado con la siguiente estructura:

\begin{lstlisting}[language=sql]
pais -> {
  id,
  name,
  slug,
  lat,
  lng,
  zoom,
  cities: [
    {
      id,
      name,
      slug,
      lat,
      lng,
      zoom,
      description,
      imageUrl,
      destinations: [
        {
          id,
          name,
          slug,
          shortDescription,
          description,
          category,
          imageUrl,
          lat,
          lng,
          cityId,
          cityName,
          activities,
          climate
        }
      ]
    }
  ]
}
\end{lstlisting}

Los filtros aplicados devuelven solo ciudades activas y destinos activos con \code{destination\_type = 'national'}. El resultado está ordenado por \code{sort\_order} y luego por nombre. Los permisos permiten acceso por \code{anon}, \code{authenticated} y \code{service\_role}.

\section{Políticas de Seguridad: Row Level Security}

RLS está habilitado en las cuatro tablas. Las políticas siguen el patrón de lectura pública y escritura exclusiva para administradores.

\begin{longtable}{L{0.25\textwidth} L{0.19\textwidth} L{0.15\textwidth} L{0.18\textwidth} L{0.13\textwidth}}
\toprule
\textbf{Política} & \textbf{Tabla} & \textbf{Operación} & \textbf{Rol} & \textbf{Condición} \\
\midrule
Public read countries & \code{countries} & SELECT & anon, authenticated & \code{is\_active = true} \\
\code{countries\_admin\_insert} & \code{countries} & INSERT & admin & Perfil con \code{role='admin'} \\
\code{countries\_admin\_update} & \code{countries} & UPDATE & admin & Perfil con \code{role='admin'} \\
\code{countries\_admin\_delete} & \code{countries} & DELETE & admin & Perfil con \code{role='admin'} \\
Public read cities & \code{cities} & SELECT & anon, authenticated & \code{is\_active = true} \\
\code{cities\_admin\_insert} & \code{cities} & INSERT & admin & Perfil con \code{role='admin'} \\
\code{cities\_admin\_update} & \code{cities} & UPDATE & admin & Perfil con \code{role='admin'} \\
\code{cities\_admin\_delete} & \code{cities} & DELETE & admin & Perfil con \code{role='admin'} \\
Public read destinations & \code{destinations} & SELECT & anon, authenticated & \code{is\_active = true} \\
\code{dest\_admin\_select} & \code{destinations} & SELECT & admin & Sin restricción adicional \\
\code{dest\_admin\_insert} & \code{destinations} & INSERT & admin & Perfil con \code{role='admin'} \\
\code{dest\_admin\_update} & \code{destinations} & UPDATE & admin & Perfil con \code{role='admin'} \\
\code{dest\_admin\_delete} & \code{destinations} & DELETE & admin & Perfil con \code{role='admin'} \\
\code{profiles\_self\_select} & \code{profiles} & SELECT & authenticated & \code{id = auth.uid()} \\
\bottomrule
\end{longtable}

\section{Extensiones PostgreSQL}

\begin{longtable}{L{0.25\textwidth} L{0.20\textwidth} L{0.45\textwidth}}
\toprule
\textbf{Extensión} & \textbf{Schema} & \textbf{Descripción} \\
\midrule
\code{uuid-ossp} & \code{extensions} & Generación de UUIDs mediante \code{uuid\_generate\_v4()}. \\
\code{pgcrypto} & \code{extensions} & Funciones criptográficas usadas internamente por Supabase Auth. \\
\code{pg\_stat\_statements} & \code{extensions} & Estadísticas de rendimiento de consultas SQL. \\
\code{supabase\_vault} & \code{vault} & Almacenamiento seguro de secretos gestionado por Supabase. \\
\bottomrule
\end{longtable}

\section{Glosario}

\begin{longtable}{L{0.26\textwidth} L{0.64\textwidth}}
\toprule
\textbf{Término} & \textbf{Definición} \\
\midrule
\code{slug} & Versión URL-amigable del nombre, en minúsculas y con guiones. Ejemplo: \code{norte-de-santander}. \\
\code{zoom} & Nivel de zoom del mapa MapLibre/Mapbox. Controla qué tan cerca o lejos aparece la cámara al seleccionar un país o ciudad. La escala va de 0, mundo, a 22, edificio. \\
\code{is\_active} & Bandera booleana que controla la visibilidad pública sin eliminar el registro. \\
\code{is\_featured} & Bandera que marca destinos para mostrar en la sección de destacados de la página de inicio. \\
\code{destination\_type} & Distingue destinos nacionales, dentro del país base de la aplicación, de destinos internacionales. \\
\code{activities} & Array de texto con actividades turísticas disponibles en el destino. \\
\code{sort\_order} & Entero que define el orden de presentación en listados. Se ordena ASC, por lo que 0 aparece primero. \\
\code{RLS} & Row Level Security: mecanismo de PostgreSQL que restringe qué filas puede ver o modificar cada usuario según políticas definidas. \\
\code{RPC} & Remote Procedure Call: función SQL invocada desde el cliente mediante la API de Supabase, por ejemplo \code{POST /rest/v1/rpc/}. \\
\bottomrule
\end{longtable}

\newpage
\section*{Referencias}
\addcontentsline{toc}{section}{Referencias}

Documentación oficial de PostgreSQL. (s. f.). \textit{PostgreSQL documentation}. Recuperado el 3 de mayo de 2026, de \url{https://www.postgresql.org/docs/}

Documentación oficial de Supabase. (s. f.). \textit{Supabase Docs}. Recuperado el 3 de mayo de 2026, de \url{https://supabase.com/docs}

Supabase. (s. f.). \textit{Row Level Security}. Recuperado el 3 de mayo de 2026, de \url{https://supabase.com/docs/guides/database/postgres/row-level-security}


\end{document}
